{\scriptsize
        \begin{tabularx}{\textwidth}{|p{0.27\columnwidth}|X|}
            \hline 
            \textbf{Notions et contenus} & \textbf{Capacités exigibles} \\
            \hline 
            \multicolumn{2}{|l|}{\textbf{6.5.3. Équation de Schrödinger dans un potentiel $V(x)$ uniforme par morceaux }} \\
            \hline
            Quantification de l'énergie dans un puits de potentiel rectangulaire de profondeur infinie.  & 
            Établir les expressions des énergies des états stationnaires.
Retrouver qualitativement l'énergie minimale à partir de l'inégalité de Heisenberg spatiale. \\
            \hline 
            Énergie de confinement quantique.   & 
            Associer le confinement d'une particule quantique à une augmentation de l'énergie cinétique. \\
        \hline    
        Évolution temporelle d'une particule confinée dans une superposition d'états. & 
        Mettre en évidence les oscillations d'une particule dont la fonction d'onde s'écrit comme la superposition de deux états stationnaires et relier la fréquence d'oscillation à la différence des énergies. \\
        \hline
        Quantification de l'énergie des états liés dans un puits de profondeur finie. Élargissement effectif du puits par les ondes évanescentes. &
        Décrire la forme des fonctions d'onde dans les différents domaines. Utiliser les conditions aux limites admises : continuité de $\phi$ et $\mathrm{d} \phi / \mathrm{dx}$. Associer la quantification de l'énergie au caractère lié de la particule.
Mener une discussion graphique. Interpréter qualitativement, à partir de l'inégalité de Heisenberg spatiale, l'abaissement des niveaux d'énergie par rapport au puits de profondeur infinie. \\
\hline
\multicolumn{2}{|l|}{\textbf{ 6.5.4. Effet tunnel }} \\
        \hline
        Effet tunnel.
Coefficient de transmission associé à une particule libre incidente sur une barrière de potentiel. & Citer quelques applications de l'effet tunnel. Définir le coefficient de transmission comme un rapport de courants de probabilités. Utiliser une expression fournie du coefficient de transmission à travers une barrière de potentiel. \\
        \hline
        \end{tabularx}
% \vspace{0.5cm}
}
